\section{Literature Review}
\label{sec:lit}

\subsection{Evolution of Predictive Modeling in Financial Markets}

Financial market forecasting has experienced a substantial development of methodology,
which has progressed over the years from primitive technical analysis to highly developed
multimodal learning architectures. The initial research was mainly based on technical
indicators and statistical methods. As an example, \citet{lee07} established that
multi-agent $Q$-learning systems could be used for daily equity trading, the limiting
aspect at that time being computational capacity and the availability of data.

The introduction of deep learning initiated a paradigm change in the process of
forecasting. Later works, including the numerical attention mechanism for dual-source
information processing proposed by \citet{liu19} and the sliding-window
metaheuristic-optimised regression for price prediction proposed by \citet{chou18}, had
significantly higher predictive accuracy. The major shortcoming of these methods,
however, was that they viewed financial markets as stationary processes, thus largely
ignoring the dynamic nature of markets in terms of regimes.

The latest level of research is marked by the creation of complex hybrid models.
\citet{nayak22} demonstrated how an enhanced chemical-reaction-optimisation algorithm can
be used in combination with dendritic neuron models, and therefore the potential of
bio-inspired computing. At the same time, \citet{farimani24} made a substantial
contribution to the sphere, developing an adaptive multimodal model for the dynamic
combination of various information sources, which is a significant step in the
comprehensive management of the intrinsic heterogeneity of financial information.
Related strands extend the same modelling apparatus to adjacent problems: forecast-horizon
model selection for commodity prices \citep{zhang20}, empirical-mode decomposition for
short-horizon series \citep{yang20}, penalised multinomial models for return direction
\citep{hu24}, and heterogeneous-source adversarial learning for financial applications
\citep{khuwaja23}.

This is described in Table~\ref{tab:evolution}, which shows the evolution from early
computational models to present-day integrative models. Every new generation is more
advanced in its capacity to process data, but one constant remains: the inability to
connect predictive indicators to practical portfolio building.

\begin{table}[!ht]
\centering
\caption{Evolution of Financial Prediction Methodologies}
\label{tab:evolution}
\footnotesize
\begin{tabular}{p{2.5cm}p{2.5cm}p{2.5cm}p{2.5cm}}
\hline
\textbf{Era} & \textbf{Primary Methods} & \textbf{Key Contributions} & \textbf{Limitations} \\
\hline
Early Computational (2004--2010) & Multi-agent systems, $Q$-learning, hybrid radial basis function (RBF) networks & \citet{lee07} -- multi-agent $Q$-learning; \citet{lee04} -- iJADE Stock Advisor with hybrid RBF & Limited data sources, computational constraints, simplistic market assumptions \\
\hline
Deep Learning Revolution (2011--2019) & Long short-term memory (LSTM) and gated recurrent unit (GRU) networks, attention mechanisms, ensemble methods & \citet{liu19} -- numerical attention methods; \citet{chou18} -- metaheuristic optimisation & Stationary market assumptions, limited regime adaptation, focus on single-asset prediction \\
\hline
Multimodal Integration (2020--Present) & Transformer architectures, multimodal fusion, LLM integration & \citet{moodi24} -- hybrid deep-learning fusion; \citet{farimani24} -- adaptive multimodal learning & Signal-allocation gap, insufficient portfolio integration, limited decision-theoretic framing \\
\hline
\end{tabular}
\end{table}

\subsection{Sentiment Analysis and Alternative Data Integration}

The use of sentiment analysis and other data streams constitutes one of the most
significant changes in modern financial forecasting. Early methodologies mostly used
lexicon-based methods and simple scoring procedures, of which the finance-specific word
lists of \citet{loughran11} are the standard reference. Modern methods, though,
increasingly exploit large language models and transformer networks to extract subtle
sentiment, as in the case of \citet{jung24}, who applied targeted masking of
numerical-change expressions and achieved significant improvements in the accuracy of
sentiment quantification.

\citet{liu24} offer a comprehensive survey of large language models in financial market
sentiment analysis, with critical analysis of the current approaches and datasets. Their
synthesis highlights the transformative ability of such models to interpret financial
discourse and, at the same time, illustrates the ongoing obstacles in terms of model
interpretability and computational requirements. Similarly, \citet{onozo24} demonstrate
that the predictive capabilities of large language models can be useful in sentiment
analysis and may also be used in the nowcasting of macroeconomic variables, which
increases the range of predictive inputs available to quantitative modellers. Within this
journal, \citet{george25} compare hybrid feature extraction across fourteen classifiers
for equity sentiment, and report that the gains depend more on the feature representation
than on the classifier.

The power of sentiment integration is still confirmed by empirical evidence.
\citet{mu23} developed a stock-price prediction model which combines investor sentiment
and optimised deep learning, and reported better performance than models based on
traditional technical indicators alone. In further support of this result,
\citet{bacco24} investigated stock prediction using LSTM networks and sentiment analysis
of tweets in a context of strong market uncertainty, which confirms the increased
informational utility of sentiment in high-volatility environments. Parallel evidence
comes from asset classes outside equities, including cryptocurrency price prediction
\citep{parekh22,otabek22}, Bitcoin range forecasting from high-dimensional indicators and
social dynamics \citep{shang25}, and electricity price prediction assisted by
sentiment agents \citep{lu25}.

A distinct and considerably older tradition measures sentiment from market data rather
than from text. \citet{bakerwurgler07} construct an investor sentiment index from
observable market quantities---trading volume, volatility and issuance activity---on the
argument that sentiment reveals itself in behaviour that can be priced, not only in
language that must first be classified. The framework proposed in this paper belongs to
this market-derived family, a point developed in Section~\ref{sec:sentiment} and validated
externally in Section~\ref{sec:sentval}.

Table~\ref{tab:sentiment_methodologies} provides an overview of the most common
sentiment-analysis strategies used in financial forecasting, tracing the history from
unsophisticated lexicon-based methods to solutions using large language models. Each
successive category shows incremental improvement in handling the complexities of
financial language, despite long-standing issues concerning computational efficiency and
transparency of the model.

\begin{table}[!ht]
\centering
\caption{Sentiment Analysis Methodologies in Financial Prediction}
\label{tab:sentiment_methodologies}
\footnotesize
\begin{tabular}{p{1.9cm}p{2.0cm}p{2.1cm}p{2.0cm}p{2.0cm}}
\hline
\textbf{Methodology Category} & \textbf{Key Studies} & \textbf{Technical Approach} & \textbf{Reported Benefits} & \textbf{Identified Gaps} \\
\hline
Lexicon-Based Approaches & \citet{loughran11} & Dictionary-based scoring, rule-based classification & Computational efficiency, interpretability & Limited context understanding, poor handling of financial jargon \\
\hline
Traditional machine learning (ML) methods & \citet{vargas22}; \citet{wang24} & Support vector machines (SVM), random forests, feature engineering & Improved accuracy over lexicons, better handling of context & Manual feature engineering, limited scalability \\
\hline
Deep Learning Architectures & \citet{mu23}; \citet{bacco24}; \citet{tai22} & LSTM, GRU and convolutional neural network (CNN) layers with embeddings & Automatic feature learning, context preservation & Computational intensity, black-box nature \\
\hline
LLM and Transformer-Based & \citet{jung24}; \citet{liu24}; \citet{onozo24} & Fine-tuned transformers, attention mechanisms, prompt engineering & Superior contextual understanding, multi-task capability & High computational cost, interpretability challenges \\
\hline
Market-Derived Measures & \citet{bakerwurgler07}; this study & Volume, volatility and participation aggregates & No text pipeline, no language dependence, fully causal & Coarser than text; cannot attribute sentiment to a named event \\
\hline
\end{tabular}
\end{table}

\subsection{Signal Fusion and Decision Integration Frameworks}

The integration of various signal typologies forms a major model in advanced financial
forecasting systems. The systematic review by \citet{zhang22} of decision fusion for
equity prediction, synthesising evidence from 76 primary studies, shows that although
fused-model architectures usually outperform unimodal ones, a high degree of diversity
remains in both fusion methods and evaluation protocols.

A range of integrative architectures has been developed in recent research.
\citet{moodi24} built a hybrid deep-learning system integrating technical indicators and
sentiment analysis, demonstrating the complementary nature of different data
representations. Further developing the field, \citet{haryono23} presented an approach
using transformer-gated recurrent units to synthesise news sentiment with technical
indicators, which emphasises the relevance of attention mechanisms in the fusion process,
and \citet{haryono24} extended the same line with aspect-based sentiment generation.
Systematic reviews of the deep-learning and technical-analysis literature report the same
pattern of accumulating architectural sophistication \citep{li20,wang20}.

Despite such methodological developments, a basic constraint persists, referred to here
as the signal-allocation gap. This weakness is a significant point of discontinuity
between sophisticated prediction models and their application to the creation of a
portfolio. As \citet{zhang22} note, the most common measure of success is still the level
of statistical forecasting accuracy, rather than economic usefulness or achieved
portfolio return. The failure is acutely apparent in the reinforcement-learning case,
such as that of \citet{kabbani22}, which largely considers single-asset trading and thus
does not address allocation at the portfolio level.

A related literature concerns how such claims should be evaluated at all.
\citet{harvey15} show that the number of specifications a researcher may examine before
reporting one inflates apparent performance to a degree that conventional significance
thresholds do not accommodate, and \citet{prado16} makes the parallel argument for
portfolio construction specifically, that a weighting scheme selected on in-sample
evidence frequently fails out of sample. Within this journal, \citet{zhangz25} report
directly that the model selected as optimal in sample is not reliably the model that
performs best out of sample. These findings motivate two design decisions taken in
Section~\ref{sec:method}: the full specification grid is reported rather than its best
cell, and a block of sessions postdating every specification decision is held out and
evaluated once.

\subsection{Market Regime Detection and Adaptation}

Recognition of market regime dependencies is an essential change in the field of
financial modelling. The foundations of regime-conscious models were formed by
\citet{wang19}, who explored economic recession prediction on the basis of multiple
behavioural factors. Their results verify that different market settings are linked with
specific behavioural signatures, and hence require distinct methods of analysis. The
econometric foundation for this view is older still: \citet{hamilton89} formalised the
notion that a single series may be generated by distinct, unobserved states with their
own dynamics, and \citet{kritzman12} carried the idea into allocation by showing that a
strategy conditioned on the inferred state can dominate its unconditional counterpart.

Regime considerations have been added at different levels in subsequent scholarship.
\citet{huang20} proposed an automated trading-point prediction framework using bicluster
mining and fuzzy inference, thereby embedding market-state dependencies implicitly in the
form of pattern recognition. Similarly, \citet{alsheebah24} combined endogenous and
exogenous variables to accommodate disparate conditions in emerging markets, though their
main aim was predictive accuracy rather than asset allocation. More recent work places
the state variable inside the allocation problem itself: \citet{li25} combine a jump model
with model predictive control, and \citet{luo26} condition both the signal and the
investable set on a dual-regime indicator. A closely related mechanism scales total risk
by recent volatility without conditioning on a discrete state at all
\citep{moreira17,hood25}; this study separates the two so that their contributions can be
compared, and Section~\ref{sec:ablation} reports the comparison.

Table~\ref{tab:paradigms} provides a comparative analysis of the portfolio construction
strategies reported in the literature, and displays how earlier methods based on
conventional optimisation techniques have been succeeded by modern multimodal and
risk-based approaches. However, in spite of the obvious rise in methodological
sophistication, substantial gaps persist, especially concerning dynamic regime adaptation
and the systematic transformation of predictive signals into the allocation choice. The
present study sits between the two responses that the literature has developed to the
instability of mean-variance inputs: robustification of the inputs themselves, as in the
multi-interval robust mean-conditional-value-at-risk formulation of
\citet{khodamoradi23}, and abandonment of the expected-return term altogether in favour
of a purely risk-based construction. It retains an expected-return term, but scales that
term explicitly so that its magnitude is commensurate with the risk term, and it separates
the composition decision from the total-risk decision so that the two can be evaluated
independently.

\begin{table}[!ht]
\centering
\caption{Portfolio Construction Methodologies in Literature}
\label{tab:paradigms}
\footnotesize
\begin{tabular}{p{2.3cm}p{2.5cm}p{2.6cm}p{2.6cm}}
\hline
\textbf{Methodology} & \textbf{Representative Studies} & \textbf{Basis of Allocation} & \textbf{Principal Limitation} \\
\hline
Traditional Optimisation & \citet{markowitz52} & Expected returns and covariance & Extreme sensitivity to estimated inputs \\
\hline
Naive Diversification & \citet{demiguel09} & None; equal weights & Ignores risk structure entirely \\
\hline
Shrinkage-Based Optimisation & \citet{ledoit04} & Regularised covariance & Still requires expected returns \\
\hline
Signal-Based Allocation & \citet{sami25} -- IntPort & Statistical averaging, forecast-based weighting & No covariance consideration, limited risk management \\
\hline
Risk-Based Construction & \citet{prado16}; \citet{asimit25} & Risk model only & Forgoes signal information by design \\
\hline
Volatility Management & \citet{moreira17}; \citet{hood25} & Total risk scaled by recent volatility & Does not address composition \\
\hline
Regime-Switching Allocation & \citet{kritzman12}; \citet{li25}; \citet{luo26} & State-dependent parameters & Component contributions rarely isolated \\
\hline
Reinforcement and Learning-Based & \citet{kabbani22}; \citet{ha26} & Learned policy or predicted returns & Single-asset focus; interpretability \\
\hline
Multimodal Fusion Approaches & \citet{moodi24}; \citet{farimani24} & Multiple signal integration, advanced ML architectures & Prediction-allocation gap, no explicit regime adaptation \\
\hline
\textbf{This study} & --- & Regime-conditional signal, plus a relative regime risk budget & Signal layer tested and found not to contribute \\
\hline
\end{tabular}
\end{table}
